\section{}

\paragraph{1137年}\eventanchor{JIN-1137-REGNAL-YEAR}{JIN-1137-REGNAL-YEAR}　金以金熙宗在位第2年作为诏令、赋役与外交文书的纪年框架。账本列出《金史》〈本纪〉、〈交聘表〉及相关志；《宋史》与《建炎以来系年要录》相关条作为来源线索；该材料可核对时间与制度称谓，但有关执行范围、群体经验、军数或后果，仍须与地方文书、物质遗存及相关政权记录互证。

\paragraph{1138年}\eventanchor{JIN-1138-REGNAL-YEAR}{JIN-1138-REGNAL-YEAR}　金以金熙宗在位第3年作为诏令、赋役与外交文书的纪年框架。账本列出《金史》〈本纪〉、〈交聘表〉及相关志；《宋史》与《建炎以来系年要录》相关条作为来源线索；该材料可核对时间与制度称谓，但有关执行范围、群体经验、军数或后果，仍须与地方文书、物质遗存及相关政权记录互证。

\paragraph{1139年}\eventanchor{JIN-1139-REGNAL-YEAR}{JIN-1139-REGNAL-YEAR}　金以金熙宗在位第4年作为诏令、赋役与外交文书的纪年框架。账本列出《金史》〈本纪〉、〈交聘表〉及相关志；《宋史》与《建炎以来系年要录》相关条作为来源线索；该材料可核对时间与制度称谓，但有关执行范围、群体经验、军数或后果，仍须与地方文书、物质遗存及相关政权记录互证。

\paragraph{1140年}\eventanchor{JIN-1140-EVENT-091}{JIN-1140-EVENT-091}　金宋在河南、淮西和陕西再度大规模作战。账本列出《金史》〈本纪〉、〈交聘表〉及相关志；《宋史》与《建炎以来系年要录》相关条作为来源线索；该材料可核对时间与制度称谓，但有关执行范围、群体经验、军数或后果，仍须与地方文书、物质遗存及相关政权记录互证。

\paragraph{1140年}\eventanchor{JIN-1140-REGNAL-YEAR}{JIN-1140-REGNAL-YEAR}　金以金熙宗在位第5年作为诏令、赋役与外交文书的纪年框架。账本列出《金史》〈本纪〉、〈交聘表〉及相关志；《宋史》与《建炎以来系年要录》相关条作为来源线索；该材料可核对时间与制度称谓，但有关执行范围、群体经验、军数或后果，仍须与地方文书、物质遗存及相关政权记录互证。

\paragraph{1141年}\eventanchor{JIN-1141-EVENT-092}{JIN-1141-EVENT-092}　金宋议定绍兴和约框架。账本列出《金史》〈本纪〉、〈交聘表〉及相关志；《宋史》与《建炎以来系年要录》相关条作为来源线索；该材料可核对时间与制度称谓，但有关执行范围、群体经验、军数或后果，仍须与地方文书、物质遗存及相关政权记录互证。

\paragraph{1141年}\eventanchor{JIN-1141-REGNAL-YEAR}{JIN-1141-REGNAL-YEAR}　金以金熙宗在位第6年作为诏令、赋役与外交文书的纪年框架。账本列出《金史》〈本纪〉、〈交聘表〉及相关志；《宋史》与《建炎以来系年要录》相关条作为来源线索；该材料可核对时间与制度称谓，但有关执行范围、群体经验、军数或后果，仍须与地方文书、物质遗存及相关政权记录互证。

\paragraph{1142年}\eventanchor{JIN-1142-REGNAL-YEAR}{JIN-1142-REGNAL-YEAR}　金以金熙宗在位第7年作为诏令、赋役与外交文书的纪年框架。账本列出《金史》〈本纪〉、〈交聘表〉及相关志；《宋史》与《建炎以来系年要录》相关条作为来源线索；该材料可核对时间与制度称谓，但有关执行范围、群体经验、军数或后果，仍须与地方文书、物质遗存及相关政权记录互证。

\paragraph{1143年}\eventanchor{JIN-1143-REGNAL-YEAR}{JIN-1143-REGNAL-YEAR}　金以金熙宗在位第8年作为诏令、赋役与外交文书的纪年框架。账本列出《金史》〈本纪〉、〈交聘表〉及相关志；《宋史》与《建炎以来系年要录》相关条作为来源线索；该材料可核对时间与制度称谓，但有关执行范围、群体经验、军数或后果，仍须与地方文书、物质遗存及相关政权记录互证。

\paragraph{1144年}\eventanchor{JIN-1144-REGNAL-YEAR}{JIN-1144-REGNAL-YEAR}　金以金熙宗在位第9年作为诏令、赋役与外交文书的纪年框架。账本列出《金史》〈本纪〉、〈交聘表〉及相关志；《宋史》与《建炎以来系年要录》相关条作为来源线索；该材料可核对时间与制度称谓，但有关执行范围、群体经验、军数或后果，仍须与地方文书、物质遗存及相关政权记录互证。

\paragraph{1145年}\eventanchor{JIN-1145-REGNAL-YEAR}{JIN-1145-REGNAL-YEAR}　金以金熙宗在位第10年作为诏令、赋役与外交文书的纪年框架。账本列出《金史》〈本纪〉、〈交聘表〉及相关志；《宋史》与《建炎以来系年要录》相关条作为来源线索；该材料可核对时间与制度称谓，但有关执行范围、群体经验、军数或后果，仍须与地方文书、物质遗存及相关政权记录互证。

\paragraph{1146年}\eventanchor{JIN-1146-REGNAL-YEAR}{JIN-1146-REGNAL-YEAR}　金以金熙宗在位第11年作为诏令、赋役与外交文书的纪年框架。账本列出《金史》〈本纪〉、〈交聘表〉及相关志；《宋史》与《建炎以来系年要录》相关条作为来源线索；该材料可核对时间与制度称谓，但有关执行范围、群体经验、军数或后果，仍须与地方文书、物质遗存及相关政权记录互证。

\paragraph{1147年}\eventanchor{JIN-1147-REGNAL-YEAR}{JIN-1147-REGNAL-YEAR}　金以金熙宗在位第12年作为诏令、赋役与外交文书的纪年框架。账本列出《金史》〈本纪〉、〈交聘表〉及相关志；《宋史》与《建炎以来系年要录》相关条作为来源线索；该材料可核对时间与制度称谓，但有关执行范围、群体经验、军数或后果，仍须与地方文书、物质遗存及相关政权记录互证。

\paragraph{1148年}\eventanchor{JIN-1148-REGNAL-YEAR}{JIN-1148-REGNAL-YEAR}　金以金熙宗在位第13年作为诏令、赋役与外交文书的纪年框架。账本列出《金史》〈本纪〉、〈交聘表〉及相关志；《宋史》与《建炎以来系年要录》相关条作为来源线索；该材料可核对时间与制度称谓，但有关执行范围、群体经验、军数或后果，仍须与地方文书、物质遗存及相关政权记录互证。

\paragraph{1149年}\eventanchor{JIN-1149-EVENT-093}{JIN-1149-EVENT-093}　完颜亮发动政变即位。账本列出《金史》〈本纪〉、〈交聘表〉及相关志；《宋史》与《建炎以来系年要录》相关条作为来源线索；该材料可核对时间与制度称谓，但有关执行范围、群体经验、军数或后果，仍须与地方文书、物质遗存及相关政权记录互证。

\paragraph{1149年}\eventanchor{JIN-1149-REGNAL-YEAR}{JIN-1149-REGNAL-YEAR}　金以金熙宗在位第14年作为诏令、赋役与外交文书的纪年框架。账本列出《金史》〈本纪〉、〈交聘表〉及相关志；《宋史》与《建炎以来系年要录》相关条作为来源线索；该材料可核对时间与制度称谓，但有关执行范围、群体经验、军数或后果，仍须与地方文书、物质遗存及相关政权记录互证。

\paragraph{1150年}\eventanchor{JIN-1150-REGNAL-YEAR}{JIN-1150-REGNAL-YEAR}　金以海陵王在位第1年作为诏令、赋役与外交文书的纪年框架。账本列出《金史》〈本纪〉、〈交聘表〉及相关志；《宋史》与《建炎以来系年要录》相关条作为来源线索；该材料可核对时间与制度称谓，但有关执行范围、群体经验、军数或后果，仍须与地方文书、物质遗存及相关政权记录互证。

\paragraph{1151年}\eventanchor{JIN-1151-REGNAL-YEAR}{JIN-1151-REGNAL-YEAR}　金以海陵王在位第2年作为诏令、赋役与外交文书的纪年框架。账本列出《金史》〈本纪〉、〈交聘表〉及相关志；《宋史》与《建炎以来系年要录》相关条作为来源线索；该材料可核对时间与制度称谓，但有关执行范围、群体经验、军数或后果，仍须与地方文书、物质遗存及相关政权记录互证。

\paragraph{1152年}\eventanchor{JIN-1152-REGNAL-YEAR}{JIN-1152-REGNAL-YEAR}　金以海陵王在位第3年作为诏令、赋役与外交文书的纪年框架。账本列出《金史》〈本纪〉、〈交聘表〉及相关志；《宋史》与《建炎以来系年要录》相关条作为来源线索；该材料可核对时间与制度称谓，但有关执行范围、群体经验、军数或后果，仍须与地方文书、物质遗存及相关政权记录互证。

\paragraph{1153年}\eventanchor{JIN-1153-EVENT-094}{JIN-1153-EVENT-094}　金正式迁都中都。账本列出《金史》〈本纪〉、〈交聘表〉及相关志；《宋史》与《建炎以来系年要录》相关条作为来源线索；该材料可核对时间与制度称谓，但有关执行范围、群体经验、军数或后果，仍须与地方文书、物质遗存及相关政权记录互证。

\paragraph{1153年}\eventanchor{JIN-1153-REGNAL-YEAR}{JIN-1153-REGNAL-YEAR}　金以海陵王在位第4年作为诏令、赋役与外交文书的纪年框架。账本列出《金史》〈本纪〉、〈交聘表〉及相关志；《宋史》与《建炎以来系年要录》相关条作为来源线索；该材料可核对时间与制度称谓，但有关执行范围、群体经验、军数或后果，仍须与地方文书、物质遗存及相关政权记录互证。

\paragraph{1154年}\eventanchor{JIN-1154-REGNAL-YEAR}{JIN-1154-REGNAL-YEAR}　金以海陵王在位第5年作为诏令、赋役与外交文书的纪年框架。账本列出《金史》〈本纪〉、〈交聘表〉及相关志；《宋史》与《建炎以来系年要录》相关条作为来源线索；该材料可核对时间与制度称谓，但有关执行范围、群体经验、军数或后果，仍须与地方文书、物质遗存及相关政权记录互证。

\paragraph{1155年}\eventanchor{JIN-1155-REGNAL-YEAR}{JIN-1155-REGNAL-YEAR}　金以海陵王在位第6年作为诏令、赋役与外交文书的纪年框架。账本列出《金史》〈本纪〉、〈交聘表〉及相关志；《宋史》与《建炎以来系年要录》相关条作为来源线索；该材料可核对时间与制度称谓，但有关执行范围、群体经验、军数或后果，仍须与地方文书、物质遗存及相关政权记录互证。

\paragraph{1156年}\eventanchor{JIN-1156-EVENT-095}{JIN-1156-EVENT-095}　完颜亮集结军队与物资准备南侵。账本列出《金史》〈本纪〉、〈交聘表〉及相关志；《宋史》与《建炎以来系年要录》相关条作为来源线索；该材料可核对时间与制度称谓，但有关执行范围、群体经验、军数或后果，仍须与地方文书、物质遗存及相关政权记录互证。

\paragraph{1156年}\eventanchor{JIN-1156-REGNAL-YEAR}{JIN-1156-REGNAL-YEAR}　金以海陵王在位第7年作为诏令、赋役与外交文书的纪年框架。账本列出《金史》〈本纪〉、〈交聘表〉及相关志；《宋史》与《建炎以来系年要录》相关条作为来源线索；该材料可核对时间与制度称谓，但有关执行范围、群体经验、军数或后果，仍须与地方文书、物质遗存及相关政权记录互证。

\paragraph{1157年}\eventanchor{JIN-1157-REGNAL-YEAR}{JIN-1157-REGNAL-YEAR}　金以海陵王在位第8年作为诏令、赋役与外交文书的纪年框架。账本列出《金史》〈本纪〉、〈交聘表〉及相关志；《宋史》与《建炎以来系年要录》相关条作为来源线索；该材料可核对时间与制度称谓，但有关执行范围、群体经验、军数或后果，仍须与地方文书、物质遗存及相关政权记录互证。

\paragraph{1158年}\eventanchor{JIN-1158-REGNAL-YEAR}{JIN-1158-REGNAL-YEAR}　金以海陵王在位第9年作为诏令、赋役与外交文书的纪年框架。账本列出《金史》〈本纪〉、〈交聘表〉及相关志；《宋史》与《建炎以来系年要录》相关条作为来源线索；该材料可核对时间与制度称谓，但有关执行范围、群体经验、军数或后果，仍须与地方文书、物质遗存及相关政权记录互证。

\paragraph{1159年}\eventanchor{JIN-1159-REGNAL-YEAR}{JIN-1159-REGNAL-YEAR}　金以海陵王在位第10年作为诏令、赋役与外交文书的纪年框架。账本列出《金史》〈本纪〉、〈交聘表〉及相关志；《宋史》与《建炎以来系年要录》相关条作为来源线索；该材料可核对时间与制度称谓，但有关执行范围、群体经验、军数或后果，仍须与地方文书、物质遗存及相关政权记录互证。

\paragraph{1160年}\eventanchor{JIN-1160-REGNAL-YEAR}{JIN-1160-REGNAL-YEAR}　金以海陵王在位第11年作为诏令、赋役与外交文书的纪年框架。账本列出《金史》〈本纪〉、〈交聘表〉及相关志；《宋史》与《建炎以来系年要录》相关条作为来源线索；该材料可核对时间与制度称谓，但有关执行范围、群体经验、军数或后果，仍须与地方文书、物质遗存及相关政权记录互证。

\paragraph{1161年}\eventanchor{JIN-1161-EVENT-096}{JIN-1161-EVENT-096}　金军南侵在采石等战场受挫，完颜亮被杀。账本列出《金史》〈本纪〉、〈交聘表〉及相关志；《宋史》与《建炎以来系年要录》相关条作为来源线索；该材料可核对时间与制度称谓，但有关执行范围、群体经验、军数或后果，仍须与地方文书、物质遗存及相关政权记录互证。

\paragraph{1161年}\eventanchor{JIN-1161-EVENT-097}{JIN-1161-EVENT-097}　完颜雍即位为金世宗。账本列出《金史》〈本纪〉、〈交聘表〉及相关志；《宋史》与《建炎以来系年要录》相关条作为来源线索；该材料可核对时间与制度称谓，但有关执行范围、群体经验、军数或后果，仍须与地方文书、物质遗存及相关政权记录互证。

\paragraph{1161年}\eventanchor{JIN-1161-REGNAL-YEAR}{JIN-1161-REGNAL-YEAR}　金以海陵王在位第12年作为诏令、赋役与外交文书的纪年框架。账本列出《金史》〈本纪〉、〈交聘表〉及相关志；《宋史》与《建炎以来系年要录》相关条作为来源线索；该材料可核对时间与制度称谓，但有关执行范围、群体经验、军数或后果，仍须与地方文书、物质遗存及相关政权记录互证。

\paragraph{1162年}\eventanchor{JIN-1162-REGNAL-YEAR}{JIN-1162-REGNAL-YEAR}　金以金世宗在位第1年作为诏令、赋役与外交文书的纪年框架。账本列出《金史》〈本纪〉、〈交聘表〉及相关志；《宋史》与《建炎以来系年要录》相关条作为来源线索；该材料可核对时间与制度称谓，但有关执行范围、群体经验、军数或后果，仍须与地方文书、物质遗存及相关政权记录互证。

\paragraph{1163年}\eventanchor{JIN-1163-REGNAL-YEAR}{JIN-1163-REGNAL-YEAR}　金以金世宗在位第2年作为诏令、赋役与外交文书的纪年框架。账本列出《金史》〈本纪〉、〈交聘表〉及相关志；《宋史》与《建炎以来系年要录》相关条作为来源线索；该材料可核对时间与制度称谓，但有关执行范围、群体经验、军数或后果，仍须与地方文书、物质遗存及相关政权记录互证。

\paragraph{1164年}\eventanchor{JIN-1164-EVENT-098}{JIN-1164-EVENT-098}　金宋缔结隆兴和议。账本列出《金史》〈本纪〉、〈交聘表〉及相关志；《宋史》与《建炎以来系年要录》相关条作为来源线索；该材料可核对时间与制度称谓，但有关执行范围、群体经验、军数或后果，仍须与地方文书、物质遗存及相关政权记录互证。

\paragraph{1164年}\eventanchor{JIN-1164-REGNAL-YEAR}{JIN-1164-REGNAL-YEAR}　金以金世宗在位第3年作为诏令、赋役与外交文书的纪年框架。账本列出《金史》〈本纪〉、〈交聘表〉及相关志；《宋史》与《建炎以来系年要录》相关条作为来源线索；该材料可核对时间与制度称谓，但有关执行范围、群体经验、军数或后果，仍须与地方文书、物质遗存及相关政权记录互证。

\paragraph{1165年}\eventanchor{JIN-1165-REGNAL-YEAR}{JIN-1165-REGNAL-YEAR}　金以金世宗在位第4年作为诏令、赋役与外交文书的纪年框架。账本列出《金史》〈本纪〉、〈交聘表〉及相关志；《宋史》与《建炎以来系年要录》相关条作为来源线索；该材料可核对时间与制度称谓，但有关执行范围、群体经验、军数或后果，仍须与地方文书、物质遗存及相关政权记录互证。

\paragraph{1166年}\eventanchor{JIN-1166-REGNAL-YEAR}{JIN-1166-REGNAL-YEAR}　金以金世宗在位第5年作为诏令、赋役与外交文书的纪年框架。账本列出《金史》〈本纪〉、〈交聘表〉及相关志；《宋史》与《建炎以来系年要录》相关条作为来源线索；该材料可核对时间与制度称谓，但有关执行范围、群体经验、军数或后果，仍须与地方文书、物质遗存及相关政权记录互证。

\paragraph{1167年}\eventanchor{JIN-1167-REGNAL-YEAR}{JIN-1167-REGNAL-YEAR}　金以金世宗在位第6年作为诏令、赋役与外交文书的纪年框架。账本列出《金史》〈本纪〉、〈交聘表〉及相关志；《宋史》与《建炎以来系年要录》相关条作为来源线索；该材料可核对时间与制度称谓，但有关执行范围、群体经验、军数或后果，仍须与地方文书、物质遗存及相关政权记录互证。

\paragraph{1168年}\eventanchor{JIN-1168-REGNAL-YEAR}{JIN-1168-REGNAL-YEAR}　金以金世宗在位第7年作为诏令、赋役与外交文书的纪年框架。账本列出《金史》〈本纪〉、〈交聘表〉及相关志；《宋史》与《建炎以来系年要录》相关条作为来源线索；该材料可核对时间与制度称谓，但有关执行范围、群体经验、军数或后果，仍须与地方文书、物质遗存及相关政权记录互证。

\paragraph{1169年}\eventanchor{JIN-1169-REGNAL-YEAR}{JIN-1169-REGNAL-YEAR}　金以金世宗在位第8年作为诏令、赋役与外交文书的纪年框架。账本列出《金史》〈本纪〉、〈交聘表〉及相关志；《宋史》与《建炎以来系年要录》相关条作为来源线索；该材料可核对时间与制度称谓，但有关执行范围、群体经验、军数或后果，仍须与地方文书、物质遗存及相关政权记录互证。

\paragraph{1170年}\eventanchor{JIN-1170-REGNAL-YEAR}{JIN-1170-REGNAL-YEAR}　金以金世宗在位第9年作为诏令、赋役与外交文书的纪年框架。账本列出《金史》〈本纪〉、〈交聘表〉及相关志；《宋史》与《建炎以来系年要录》相关条作为来源线索；该材料可核对时间与制度称谓，但有关执行范围、群体经验、军数或后果，仍须与地方文书、物质遗存及相关政权记录互证。

\paragraph{1171年}\eventanchor{JIN-1171-REGNAL-YEAR}{JIN-1171-REGNAL-YEAR}　金以金世宗在位第10年作为诏令、赋役与外交文书的纪年框架。账本列出《金史》〈本纪〉、〈交聘表〉及相关志；《宋史》与《建炎以来系年要录》相关条作为来源线索；该材料可核对时间与制度称谓，但有关执行范围、群体经验、军数或后果，仍须与地方文书、物质遗存及相关政权记录互证。

\paragraph{1172年}\eventanchor{JIN-1172-REGNAL-YEAR}{JIN-1172-REGNAL-YEAR}　金以金世宗在位第11年作为诏令、赋役与外交文书的纪年框架。账本列出《金史》〈本纪〉、〈交聘表〉及相关志；《宋史》与《建炎以来系年要录》相关条作为来源线索；该材料可核对时间与制度称谓，但有关执行范围、群体经验、军数或后果，仍须与地方文书、物质遗存及相关政权记录互证。

\paragraph{1173年}\eventanchor{JIN-1173-REGNAL-YEAR}{JIN-1173-REGNAL-YEAR}　金以金世宗在位第12年作为诏令、赋役与外交文书的纪年框架。账本列出《金史》〈本纪〉、〈交聘表〉及相关志；《宋史》与《建炎以来系年要录》相关条作为来源线索；该材料可核对时间与制度称谓，但有关执行范围、群体经验、军数或后果，仍须与地方文书、物质遗存及相关政权记录互证。

\paragraph{1174年}\eventanchor{JIN-1174-REGNAL-YEAR}{JIN-1174-REGNAL-YEAR}　金以金世宗在位第13年作为诏令、赋役与外交文书的纪年框架。账本列出《金史》〈本纪〉、〈交聘表〉及相关志；《宋史》与《建炎以来系年要录》相关条作为来源线索；该材料可核对时间与制度称谓，但有关执行范围、群体经验、军数或后果，仍须与地方文书、物质遗存及相关政权记录互证。

\paragraph{1175年}\eventanchor{JIN-1175-REGNAL-YEAR}{JIN-1175-REGNAL-YEAR}　金以金世宗在位第14年作为诏令、赋役与外交文书的纪年框架。账本列出《金史》〈本纪〉、〈交聘表〉及相关志；《宋史》与《建炎以来系年要录》相关条作为来源线索；该材料可核对时间与制度称谓，但有关执行范围、群体经验、军数或后果，仍须与地方文书、物质遗存及相关政权记录互证。

\paragraph{1176年}\eventanchor{JIN-1176-REGNAL-YEAR}{JIN-1176-REGNAL-YEAR}　金以金世宗在位第15年作为诏令、赋役与外交文书的纪年框架。账本列出《金史》〈本纪〉、〈交聘表〉及相关志；《宋史》与《建炎以来系年要录》相关条作为来源线索；该材料可核对时间与制度称谓，但有关执行范围、群体经验、军数或后果，仍须与地方文书、物质遗存及相关政权记录互证。

\paragraph{1177年}\eventanchor{JIN-1177-REGNAL-YEAR}{JIN-1177-REGNAL-YEAR}　金以金世宗在位第16年作为诏令、赋役与外交文书的纪年框架。账本列出《金史》〈本纪〉、〈交聘表〉及相关志；《宋史》与《建炎以来系年要录》相关条作为来源线索；该材料可核对时间与制度称谓，但有关执行范围、群体经验、军数或后果，仍须与地方文书、物质遗存及相关政权记录互证。

\paragraph{1178年}\eventanchor{JIN-1178-REGNAL-YEAR}{JIN-1178-REGNAL-YEAR}　金以金世宗在位第17年作为诏令、赋役与外交文书的纪年框架。账本列出《金史》〈本纪〉、〈交聘表〉及相关志；《宋史》与《建炎以来系年要录》相关条作为来源线索；该材料可核对时间与制度称谓，但有关执行范围、群体经验、军数或后果，仍须与地方文书、物质遗存及相关政权记录互证。

